\section{La inteligencia no es una línea de puntaje en lenguaje}

Al final de la entrevista, Xie Saining (谢赛宁) sostiene que la AGI es un falso problema y habla de la inteligencia animal y de la arrogancia humana. Cita libros sobre cognición animal y señala que cada ser vivo tiene sus propias formas de percibir y de actuar: el razonamiento de los chimpancés, el almacenamiento de alimento de las aves, la comunicación de las ballenas, el olfato de los perros, el oído de los murciélagos. La implicación técnica es que la inteligencia no se ordena sobre una línea de puntaje de un examen de lenguaje, sino que está ligada al cuerpo, al entorno, a los objetivos, a los canales de percepción y a las tareas de supervivencia.

\begin{figure}[H]
\centering
\includegraphics[width=0.94\textwidth]{intelligence-spectrum.png}
\caption{La inteligencia no es una línea de puntaje en lenguaje: del texto y el código a la visión, la acción, el cuerpo y el entorno, y los objetivos sociales.}
\end{figure}

\begin{knowledgebox}{Lectura de la figura: por qué importa una perspectiva no antropocéntrica}
La figura descompone la inteligencia en varias dimensiones: texto/código, visión/espacio, acción/control, cuerpo/entorno y sociedad/objetivos. Que un LLM rinda bien en texto y código no significa que posea automáticamente capacidades en el plano del cuerpo, del entorno y de la acción. El sentido de los modelos del mundo es precisamente devolver la inteligencia a la percepción, la predicción, la acción y la retroalimentación.
\end{knowledgebox}

\subsection{La inteligencia de una ardilla y las tareas domésticas de un niño}

En la entrevista se cita una perspectiva de Rich Sutton: en lugar de maravillarse porque un modelo escribe código o gana medallas en competencias de matemáticas, conviene pensar en lo difícil que es construir una ardilla capaz de sobrevivir en el mundo real. La afirmación parece exagerada, pero desplaza la pregunta de la «capacidad para aprobar exámenes humanos» a la capacidad básica de «un ser vivo que actúa en el mundo».

Xie Saining añade que un niño de pocos años, o uno de doce, puede hacer muchas tareas domésticas que los robots actuales todavía no realizan de forma confiable. No se trata de menospreciar a los LLM, sino de recordar que la percepción, el movimiento, la recuperación ante errores, los objetivos y el sentido común en el mundo real siguen siendo grandes problemas sin resolver.

\begin{warningbox}{El lema de la AGI tiende a ocultar las capacidades concretas}
Si no se definen el entorno, el cuerpo, los objetivos y la evaluación, la AGI se convierte con facilidad en una etiqueta vacía. Una pregunta más operable es: ¿puede el sistema, en un entorno real determinado, percibir, predecir y actuar de forma confiable, recuperarse después de equivocarse y transferir gradualmente lo aprendido a tareas nuevas?
\end{warningbox}

\subsection{Resumen de la sección}

La inteligencia no tiene una sola escala. La capacidad lingüística es importante, pero la inteligencia en el mundo real también abarca la visión, el cuerpo, la acción, la retroalimentación y los objetivos sociales. La línea de investigación de los modelos del mundo busca cubrir justamente estas dimensiones que el puntaje en lenguaje deja ocultas.
